\chapter{Código Fuente del Sistema}
\label{apx:codigo}

\par En este apartado se recopila el código fuente completo de los módulos desarrollados para el sistema de monitoreo remoto de equipos fiscales. El material se organiza en cuatro bloques: el firmware del Módulo de Transmisión IoT (MT-IoT) en lenguaje C, las estructuras de mensajería en formato JSON, las herramientas de simulación y el \textit{worker} de servidor en Python, y la capa de interfaz reactiva en HTML. Cada bloque conserva los comentarios originales que facilitan su interpretación por terceros.

\section{Firmware del Módulo de Transmisión IoT (MT-IoT)}

\subsection{Sanitización y Parseo de la Trama SV2}

\begin{lstlisting}[language=C, caption={Mitigación del error Off-By-One durante el parseo de la trama SV2.}, label={code:offbyone}, frame=single, basicstyle=\footnotesize\ttfamily, keywordstyle=\bfseries]
int bytes_leidos = receive_response(rx_buffer, sizeof(rx_buffer), 1000);

if (bytes_leidos > 15 && validate_frame(rx_buffer, bytes_leidos)) {
    int etx_pos = bytes_leidos - 2; // Ajuste para esquivar salto de linea
    int fw_len = 10; // Ejemplo: 020507GD00
    
    if (etx_pos >= fw_len + 1) {
        int start_idx = etx_pos - fw_len - 1;
        if (xSemaphoreTake(mqtt_data_mutex, pdMS_TO_TICKS(100)) == pdTRUE) {
            memset(mqtt_data.fw_printer, 0, sizeof(mqtt_data.fw_printer));
            // Copiado seguro hacia la estructura
            strncpy(mqtt_data.fw_printer, (char*)&rx_buffer[start_idx], fw_len);
            mqtt_data.fw_printer[fw_len] = '\0'; // Terminador nulo manual
            
            ESP_LOGI(TAG_PRINTER, "Version detectada via SV2: %s", 
                     mqtt_data.fw_printer);
            xSemaphoreGive(mqtt_data_mutex);
        }
        return true;
    }
}
\end{lstlisting}

\subsection{Interrogación Serial No Bloqueante}

\begin{lstlisting}[language=C, caption={Lógica no bloqueante de la tarea de interrogación serial (Polling).}, label={code:task_delay}, frame=single, basicstyle=\footnotesize\ttfamily, keywordstyle=\bfseries]
void printer_polling_task(void *pvParameters) {
    uint8_t enq_cmd = 0x05; // Caracter ENQ segun el manual
    
    while(1) {
        // Transmision asincrona del byte de estado
        uart_write_bytes(UART_NUM_1, (const char*)&enq_cmd, 1);
        
        // El hilo cede voluntariamente el control al Scheduler por 5 segundos
        vTaskDelay(pdMS_TO_TICKS(5000)); 
    }
}
\end{lstlisting}

\subsection{Validación de Integridad de Trama}

\begin{lstlisting}[language=C, caption={Lógica de validación de trama e intercepción de ruido electromagnético.}, label={code:validate_frame}]
// Validar integridad estructural y matematica (LRC) de la trama
if (!validate_frame(rx_buffer, received)) {
    ESP_LOGW(TAG_PRINTER, "Trama invalida detectada por ruido/Checksum");
    printer_handle_error(ERROR_INVALID_RESPONSE);
    // Liberacion del semaforo para evitar bloqueos del sistema
    xSemaphoreGive(spi_bus_mutex);
    continue;
}
\end{lstlisting}

\subsection{Exclusión Mutua de Memoria y Bus Físico}

\begin{lstlisting}[language=C, caption={Uso de doble exclusión mutua para proteger RAM y hardware físico.}, label={code:mutex_bus}]
// 1. Tomar el control exclusivo del cable fisico sincrono
if (xSemaphoreTake(bus_local_mutex, portMAX_DELAY) == pdTRUE) {
    
    // Validar integridad estructural y matematica de la trama
    if (validate_frame(rx_buffer, received)) {
        
        // 2. Tomar el control de la memoria RAM para escritura
        if (xSemaphoreTake(mqtt_data_mutex, pdMS_TO_TICKS(100)) == pdTRUE) {
            strncpy(mqtt_data.fw_printer, (char*)&rx_buffer[start_idx], fw_len);
            mqtt_data.fw_printer[fw_len] = '\0';
            xSemaphoreGive(mqtt_data_mutex); // Liberar RAM
        }
    }
    xSemaphoreGive(bus_local_mutex); // Liberar cable fisico
}
\end{lstlisting}

\subsection{Resiliencia y Reconexión MQTT}

\begin{lstlisting}[language=C, caption={Lógica de resiliencia y reconexión en la tarea MQTT.}, label={code:mqtt_reconnect}]
// VERIFICACION DE CONEXION DENTRO DE LA TAREA ASINCRONA
if (!mqtt_connected) {
    LOG_W(TAG_MQTT, "MQTT desconectado, reintentando en breve...");
    // Delay de cortesia para no saturar la red (Retardo exponencial)
    vTaskDelay(pdMS_TO_TICKS(RECONNECT_DELAY_MS)); 
    continue;
}
\end{lstlisting}

\subsection{Configuración Segura del Cliente MQTT}

\begin{lstlisting}[language=C, caption={Configuración de inicialización del cliente MQTT con seguridad TLS estricta.}, label={code:mqtt_tls}]
esp_mqtt_client_config_t mqtt_cfg = {
    .broker.address.hostname = "832b8689599f4045be005c116bc416f0.s1.eu.hivemq.cloud",
    .broker.address.port = 8883,
    .broker.address.transport = MQTT_TRANSPORT_OVER_SSL,
    .broker.verification.crt_bundle_attach = esp_crt_bundle_attach,
    // ... configuracion de credenciales y sesion ...
};
\end{lstlisting}

\subsection{Inyección de Metadatos mediante User Properties}

\begin{lstlisting}[language=C, caption={Inyección dinámica de metadatos mediante User Properties de MQTT v5.0.}, label={code:user_props}]
esp_mqtt5_user_property_item_t user_prop_items[] = {
    {"modelo_dispositivo", "HKA80"},
    {"firmware_version", ISMART_VERSION},
    {"serial", serial_seguro},
    {"tipo_dato", tipo_dato},
    {"wifi_rssi", rssi_str} // Capturado via esp_wifi_sta_get_ap_info
};

mqtt5_user_property_handle_t user_prop_handle = NULL;
esp_mqtt5_client_set_user_property(&user_prop_handle, user_prop_items, 5);
\end{lstlisting}

\subsection{Testamento (Last Will and Testament)}

\begin{lstlisting}[language=json, caption={Mensaje de última voluntad (\textit{Last Will and Testament}) registrado en el broker.}, label={code:lwt_payload}, frame=single, basicstyle=\footnotesize\ttfamily]
{"st": "offline", "alerta": "conexion_perdida_abruptamente"}
\end{lstlisting}

\subsection{Reanudación de Descargas mediante HTTP Range}

\begin{lstlisting}[language=C, caption={Inyección de cabecera HTTP Range para reanudación de descargas interrumpidas.}, label={code:http_range}]
if (total_bytes_recibidos > 0) {
    char range_header[32];
    // Solicitamos al servidor continuar desde el byte de corte exacto
    snprintf(range_header, sizeof(range_header), "bytes=%d-", total_bytes_recibidos);
    esp_http_client_set_header(client, "Range", range_header);
}
esp_err_t err = esp_http_client_open(client, 0);
\end{lstlisting}

\subsection{Empaquetado de Firmware Segmentado hacia el Periférico}

\begin{lstlisting}[language=C, caption={Empaquetado y transmisión de firmware segmentado hacia el periférico local.}, label={code:local_ota}]
// Empaquetado estricto de capa fisica
transaccionar_byte_bus(CMD_PACKET); // Comando de insercion
enviar_chunk_bus((uint8_t*)&num_pack, 4); // Secuencial del paquete
enviar_chunk_bus((uint8_t*)&total_packs, 4); // Total esperado
uint16_t tam_actual = (uint16_t)len;
enviar_chunk_bus((uint8_t*)&tam_actual, 2); // Longitud del fragmento
enviar_chunk_bus(data, len); // Bloque de datos crudos
enviar_chunk_bus((uint8_t*)&checksum_paquete, 4); // Redundancia

if (transaccionar_byte_bus(0x00) == ACK_BYTE) {
    return ESP_OK; // Bloque integrado exitosamente
}
\end{lstlisting}

\section{Estructuras de Mensajería (JSON)}

\subsection{Telemetría del Status S1}

\begin{lstlisting}[language=json, caption={Estructura JSON del Status S1 publicado por el MT-IoT hacia el broker MQTT.}, label={code:json_mqtt}, frame=single, basicstyle=\footnotesize\ttfamily, keywordstyle=\bfseries]
{
  "timestamp": 164370,
  "sts1": "0x60",
  "sts2": "0x40",
  "estado": "Modo Fiscal y en Espera",
  "error": "Ningun error",
  "cajero": "0010",
  "ventas": "00000000000000000",
  "fact_emitidas": "00000",
  "rif": "J3121711970",
  "registro": "Z1F9999988"
}
\end{lstlisting}

\subsection{Comando de Actualización Remota}

\begin{lstlisting}[language=json, caption={Estructura del payload MQTT para iniciar un proceso de actualización remota.}, label={code:json_ota}, frame=single, basicstyle=\footnotesize\ttfamily, keywordstyle=\bfseries]
{
  "comando": "INICIAR_OTA",
  "objetivo": "ismart",
  "version": "v020100",
  "url_descarga": "https://servidor-seguro.com/media/firmwares/update.bin"
}
\end{lstlisting}

\section{Herramientas de Simulación}

\begin{lstlisting}[language=Python, caption={Simulador fiscal con cálculo de redundancia (LRC) y manejo de estados.}, label={code:simulador_python}, frame=single, basicstyle=\footnotesize\ttfamily, keywordstyle=\bfseries]
def procesar_trama_completa(ser, trama):
    # Extraccion de datos y byte de comprobacion LRC
    data_etx = trama[1:-1] 
    lrc_rx = trama[-1] 
    lrc_calc = calcular_lrc(data_etx) 
    
    if lrc_rx == lrc_calc: 
        ser.write(bytes([0x06])) # Enviar ACK (0x06)
        comando = data_etx[:-1].decode('utf-8', errors='ignore')
        
        if comando == "S1": 
            enviar_respuesta_s1(ser) # Trama con variables fiscales
        elif comando == "SV2": 
            enviar_respuesta_sv2(ser) # Version de firmware
    else: 
        ser.write(bytes([0x15])) # Enviar NAK por error de LRC
\end{lstlisting}

\section{Worker de Telemetría en el Servidor}

\begin{lstlisting}[language=Python, caption={Lógica de suscripción y manejo de desconexiones LWT en el Worker MQTT.}, label={code:listen_mqtt}, frame=single, basicstyle=\footnotesize\ttfamily, keywordstyle=\bfseries]
def on_message(client, userdata, msg):
    partes_topico = msg.topic.split('/')
    serial = partes_topico[2]
    tipo = partes_topico[3] # ej. 'alertas_conexion' o 'json_completo'
    
    payload = json.loads(msg.payload.decode('utf-8'))
    
    # Interceptacion de Alertas de Red (Testamento LWT)
    if tipo == "alertas_conexion" and payload.get("st") == "offline":
        # Clonacion transaccional del ultimo log valido para no perder el estado
        dispositivo = Dispositivo.objects.get(serial=serial)
        LogDispositivo.objects.create(
            dispositivo=dispositivo,
            evento="Desconexion Abrupta (LWT)",
            detalles=payload
        )
\end{lstlisting}

\section{Interfaz Reactiva Web}

\begin{lstlisting}[language=HTML, caption={Captura reactiva de eventos MQTT mediante atributos declarativos de HTMX.}, label={code:htmx_log}, frame=single, basicstyle=\footnotesize\ttfamily, keywordstyle=\bfseries]
<!-- log_capture.html -->
<!-- hx-trigger escucha el evento despachado por el cliente WebSockets (MQTT.js) -->
<div id="vista-log-capture" class="p-4"
    hx-get="/log_capture/?buscar={{ request.GET.buscar }}"
    hx-trigger="mqttNuevoLog from:body"
    hx-target="#vista-log-capture"
    hx-swap="outerHTML">
    
    <div class="d-flex justify-content-between align-items-center mb-4">
        <h4>Historial de Auditoria por Dispositivo</h4>
        <small class="text-muted">Conservacion automatica: 90 dias</small>
    </div>
    <!-- ... renderizado dinamico de la tabla ... -->
</div>
\end{lstlisting}
